%% ch01.tex --- What Machine Learning Is
%% Numbers in Section 1.8 are computed by verify/ch01_flat_prices.py.

\chapter{What Machine Learning Is}
\label{ch:ch01}

\epigraph{Programming computers to learn from experience should eventually
eliminate the need for much of this detailed programming effort.}%
{---\ Arthur L.\ Samuel, 1959}

\begin{chapterlearning}
After working through this chapter you will be able to:
\begin{enumerate}
  \item Say what distinguishes a machine learning solution from an ordinary
        program, and recognise the situations in which the distinction matters.
  \item State Mitchell's definition of learning and instantiate
        $\langle T, P, E\rangle$ for a problem you are given.
  \item Name the four components present in every supervised learning
        algorithm: the data, the hypothesis space, the cost, and the optimizer.
  \item Classify a problem as supervised regression, supervised
        classification, unsupervised, or reinforcement learning, and justify
        the classification.
  \item Read a scatter plot with a fitted model on it, and compute the cost of
        that model on the plotted data.
  \item Explain why a model must be judged against a baseline and on data it
        has not seen.
  \item Identify problems for which machine learning is the wrong tool.
\end{enumerate}
\end{chapterlearning}

\section{Two Ways to Make a Computer Do Something}
\label{sec:two-ways}

Consider two tasks. The first: given a list of one million numbers, sort it into
ascending order. The second: given a photograph, decide whether it contains a
cat.

You already know how to do the first. You choose an algorithm, you write it
down, you can prove it correct, and you can state exactly how long it will take.
The rules that solve the problem are known in advance, and programming consists
of writing them out in a language a machine can execute.

Now try to do the same for the second task. Write down the rules that decide
whether a photograph contains a cat. You will find yourself reaching for
conditions like \emph{two triangular ears near the top} or \emph{a pair of
elliptical highlights where the eyes are}, and every one of them will fail on
photographs a three-year-old classifies instantly: a cat seen from behind, a
black cat in shadow, a cat behind a chair leg. The rules exist --- you apply them
several times a day --- but nobody can state them.

This is the gap machine learning fills. Instead of supplying the rules, you
supply examples of the correct answers, and a learning algorithm searches for
rules consistent with them. \Cref{fig:two-ways} draws the difference.

\begin{figure}[htbp]
  \centering
  \begin{minipage}[t]{0.47\linewidth}
    \centering
    {\small\textbf{Classical programming}\par}
    \vspace{4pt}
    \begin{tikzpicture}[node distance=5mm,
                        every node/.style={font=\scriptsize}]
      \node[draw,rounded corners,fill=MLLight,minimum width=1.5cm] (r) {Rules};
      \node[draw,rounded corners,fill=MLLight,minimum width=1.5cm,below=of r] (d) {Data};
      \node[draw,fill=MLBlue,text=white,rounded corners,
            right=7mm of $(r)!0.5!(d)$] (p) {Program};
      \node[draw,rounded corners,fill=MLAmber!25,right=7mm of p] (a) {Answers};
      \draw[-{Stealth}] (r) -- (p);
      \draw[-{Stealth}] (d) -- (p);
      \draw[-{Stealth}] (p) -- (a);
    \end{tikzpicture}
    \vspace{4pt}
    {\scriptsize The human supplies the rules.\par}
  \end{minipage}\hfill
  \begin{minipage}[t]{0.47\linewidth}
    \centering
    {\small\textbf{Machine learning}\par}
    \vspace{4pt}
    \begin{tikzpicture}[node distance=5mm,
                        every node/.style={font=\scriptsize}]
      \node[draw,rounded corners,fill=MLLight,minimum width=1.5cm] (d) {Data};
      \node[draw,rounded corners,fill=MLLight,minimum width=1.5cm,below=of d] (a) {Answers};
      \node[draw,fill=MLTeal,text=white,rounded corners,
            right=7mm of $(d)!0.5!(a)$] (t) {Training};
      \node[draw,rounded corners,fill=MLAmber!25,right=7mm of t] (r) {Rules};
      \draw[-{Stealth}] (d) -- (t);
      \draw[-{Stealth}] (a) -- (t);
      \draw[-{Stealth}] (t) -- (r);
    \end{tikzpicture}
    \vspace{4pt}
    {\scriptsize The machine infers the rules from examples.\par}
  \end{minipage}
  \caption{The two arrangements. Machine learning inverts the usual flow: rules
  become an output rather than an input.}
  \label{fig:two-ways}
\end{figure}

The inversion has a consequence that is easy to state and hard to internalise.
In classical programming, the program is the thing you write and the answers are
what you get. In machine learning, the answers are what you supply and the
program is what you get. Everything that is difficult about the subject follows
from this: the quality of what you get out is bounded by the quality of the
examples you put in, and the rules you get are not written in a form you can
easily read.

\begin{keyinsight}{When to reach for machine learning}
Machine learning is not magic, and it is not a replacement for algorithms. It is
the tool of choice in exactly one situation: when a rule exists but cannot be
written down, and examples of the rule's output are available in quantity. If
you can state the rule, state it --- the resulting program will be faster,
smaller, testable, and explainable.
\end{keyinsight}

\section{Where Machine Learning Is Used}
\label{sec:applications}

You have already used a dozen learned models today. \Cref{tab:applications}
lists a few, with the part of each system that is learned rather than
programmed.

\begin{table}[htbp]
  \centering
  \small
  \renewcommand{\arraystretch}{1.25}
  \begin{tabular}{@{}L{3.3cm}L{3.9cm}L{4.0cm}@{}}
    \toprule
    \textbf{System} & \textbf{What is learned} & \textbf{From what examples}\\
    \midrule
    Web search        & how well a page answers a query & billions of past
                        queries and the results people clicked\\
    Spam filtering    & whether a message is unwanted    & messages users marked
                        as spam or not\\
    Speech recognition& which words an audio clip contains & transcribed audio\\
    Machine translation & a sentence in another language & aligned sentence
                        pairs\\
    Medical imaging   & whether a scan is suspicious     & scans labelled by
                        radiologists\\
    Credit scoring    & the probability a loan is repaid & past loans and their
                        outcomes\\
    Predictive maintenance & when a machine will fail    & sensor logs with
                        recorded failures\\
    Recommendation    & what a user will want next       & what similar users
                        chose\\
    \bottomrule
  \end{tabular}
  \caption{Applications, and the component in each that is learned from data.}
  \label{tab:applications}
\end{table}

The common thread is worth naming, because it is also the checklist you will use
when someone proposes a machine learning project. In every row: a large pool of
\emph{recorded past decisions} exists; the decision rule is too intricate,
too varied, or too personal to hand-code; and a small error rate is tolerable.
Take away any one of the three and the case for learning weakens. There is no
pool of recorded decisions for a system nobody has operated yet. The rule for
computing sales tax is intricate but perfectly writable, so learning it from
examples of receipts would be an act of self-harm. And no error rate is
tolerable in the code that computes a bank balance.

\section{Defining Learning}
\label{sec:defining}

Two definitions have stuck, and they do different jobs.

\begin{defbox}{Machine learning, informally}
Machine learning is the field of study that gives computers the ability to learn
without being explicitly programmed.
\end{defbox}

This is the popular paraphrase of Arthur Samuel's position, and it is a good
slogan: it says what the field is \emph{for}. What it does not do is tell you
whether a particular system is working. For that we need a definition with
moving parts.

\begin{defbox}{Machine learning, operationally (Mitchell)}
A computer program is said to \term{learn} from experience $E$ with respect to
some class of tasks $T$ and performance measure $P$, if its performance at tasks
in $T$, as measured by $P$, improves with experience $E$.
\end{defbox}

Mitchell's definition earns its place by forcing you to name three things before
you write any code \parencite{mitchell1997}. Name the task $T$, and you have
said what the system does. Name the performance measure $P$, and you have said
how you will know whether it works. Name the experience $E$, and you have said
what data you need. A project that cannot fill in all three boxes is not yet a
project. In my experience the box most often left empty is $P$, and the projects
that leave it empty are the ones that run for months without anyone being able
to say whether they are succeeding.

\begin{table}[htbp]
  \centering
  \small
  \renewcommand{\arraystretch}{1.3}
  \begin{tabular}{@{}L{2.3cm}L{2.9cm}L{3.1cm}L{2.9cm}@{}}
    \toprule
    \textbf{Problem} & \textbf{Task $T$} & \textbf{Performance $P$} &
    \textbf{Experience $E$}\\
    \midrule
    Checkers player & play a game of checkers &
      fraction of games won against opponents &
      games played against itself\\
    Spam filter & label a message spam or not &
      fraction of messages labelled correctly &
      a corpus of user-flagged messages\\
    Price estimator & predict a sale price from floor area &
      mean squared prediction error &
      records of past sales\\
    Handwriting reader & identify the digit in an image &
      fraction of digits read correctly &
      images labelled by hand\\
    \bottomrule
  \end{tabular}
  \caption{Four problems with $\langle T,P,E\rangle$ filled in. Doing this
  before writing code is the cheapest useful step in a project.}
  \label{tab:tpe}
\end{table}

\begin{pitfall}{Accuracy is not automatically the right $P$}
Suppose a spam filter is evaluated on a mailbox in which $99\%$ of messages are
legitimate. The program \code{return "not spam"} --- which contains no learning,
no data, and no intelligence --- achieves $99\%$ accuracy. Any real filter that
scores $97\%$ looks worse by this measure while being infinitely more useful.
Choosing $P$ is a design decision with consequences, and on skewed data,
accuracy is usually the wrong choice. We return to this in \cref{ch:ch11} with
measures that survive skew.
\end{pitfall}

\begin{researchspotlight}{Samuel (1959) --- Studies in Machine Learning Using the Game of Checkers}
\textbf{The problem.} In the 1950s, getting a computer to play a board game well
meant a programmer encoding their own judgement of a good position. Samuel, at
IBM, wanted the machine to acquire that judgement itself.

\textbf{The idea.} His program evaluated a board position with a weighted sum of
simple features --- piece count, kings, control of the centre. The weights were
not set by hand. The program played, recorded which positions preceded wins, and
adjusted the weights toward evaluations that had proved predictive: an early
form of what \cref{ch:ch17} calls temporal-difference learning. A second scheme,
rote learning, cached previously evaluated positions.

\textbf{The result.} The program came to play a better game than Samuel himself,
who was a weak player, and it did so on an IBM 704 with a few kilobytes of
memory. That a program could exceed its author's ability at the very task it was
written for was the paper's real contribution, and it remains the clearest
one-sentence argument for the field.

\textbf{Reading it today.} The famous definition quoted at the top of this
section is a paraphrase; the phrase ``without being explicitly programmed'' does
not appear verbatim in the paper. What does appear, in the first paragraph, is
the epigraph to this chapter --- a more modest and more accurate claim.
\parencite{samuel1959}
\end{researchspotlight}

\section{The Four Parts of a Learning Algorithm}
\label{sec:four-parts}

Every supervised learning method in this book --- linear regression, logistic
regression, neural networks, decision trees --- is assembled from the same four
components. Learning to see them separately is the single most useful habit you
can form in this chapter, because when a model misbehaves, the fix always lies
in exactly one of the four.

\subsection{The data}

\begin{defbox}{Training set}
A \term{training set} is a collection of $m$ examples
\[
  \Dset \;=\; \bigl\{(\xii{1},\yii{1}),\,(\xii{2},\yii{2}),\,\dots,\,
                     (\xii{m},\yii{m})\bigr\},
  \quad \xii{i}\in\Xspace,\; \yii{i}\in\Yspace,
\]
where $\xii{i}$ is the \term{input} (or \term{feature}) of the $i$-th example and
$\yii{i}$ is its \term{target} (or \term{label}). The superscript in parentheses
is an index, never an exponent.
\end{defbox}

That last sentence is not pedantry. Throughout this book $\xii{2}$ means the
second training example and $x^2$ means $x$ squared, and the two appear in the
same equation often enough that confusing them will cost you an exercise.

\subsection{The hypothesis space}

The learning algorithm does not search all possible functions --- there are far
too many, and most of them merely memorise the training data. Instead you choose,
in advance, a family of candidate functions.

\begin{defbox}{Hypothesis space}
The \term{hypothesis space} $\Hspace$ is the set of functions
$f:\Xspace\to\Yspace$ that the learning algorithm is permitted to return. It is
chosen by the designer, before training, and it is usually described by a small
number of \term{parameters}.
\end{defbox}

For example, taking $\Hspace$ to be the set of straight lines
$\{f(x)=wx+b : w,b\in\R\}$ gives \term{linear regression}, the subject of
\cref{ch:ch02}. Every choice of $(w,b)$ is one hypothesis; learning means
choosing among them. Choosing $\Hspace$ well is the central modelling decision:
too small and no member of it fits the data (\cref{ch:ch10} calls this
\emph{bias}), too large and many members fit the training data while disagreeing
wildly elsewhere (\emph{variance}, and \cref{ch:ch07}'s \emph{overfitting}).

\subsection{The cost function}

Given a candidate $f$, we need a number saying how badly it does on the training
set. That number is the \term{cost} (or objective) $J$. For regression, the usual
choice measures squared error,
\[
  J \;=\; \frac{1}{2m}\sum_{i=1}^{m}\bigl(f(\xii{i}) - \yii{i}\bigr)^2 ,
\]
and \cref{ch:ch02} explains every part of it, including the otherwise mysterious
$\tfrac12$. For now, note only that $J$ is a function of the parameters, not of
the data: the data is fixed, and the parameters are what we vary.

\subsection{The optimizer}

Finally, something must search $\Hspace$ for the hypothesis of least cost:
\[
  f^\star \;=\; \argmin_{f\in\Hspace} J(f).
\]
Sometimes this can be solved with algebra in one step (\cref{ch:ch04}'s normal
equation). Usually it cannot, and we use \term{gradient descent}
(\cref{ch:ch03}), which starts from a guess and repeatedly steps downhill.

\begin{keyinsight}{Learning is search}
Learning is a search over a hypothesis space, guided by a cost function. The
model you get is not ``what the data says''; it is the cheapest member of a
family \emph{you} chose, measured by a criterion \emph{you} chose. Change the
family or the criterion and you get a different answer from identical data.
\end{keyinsight}

\begin{figure}[htbp]
  \centering
  \begin{tikzpicture}[every node/.style={font=\scriptsize},node distance=7mm]
    \node[draw,rounded corners,fill=MLLight,align=center,minimum height=9mm]
      (D) {Training set\\ $(\xii{i},\yii{i})_{i=1}^{m}$};
    \node[draw,fill=MLTeal,text=white,rounded corners,right=of D,
          minimum height=9mm,align=center] (L) {Learning algorithm\\
          \scriptsize ($\Hspace$, $J$, optimizer)};
    \node[draw,fill=MLBlue,text=white,rounded corners,right=of L,
          minimum height=9mm] (f) {model $f$};
    \node[right=12mm of f,align=center] (x) {new input $x$};
    \node[below=9mm of f,align=center] (yh) {prediction $\yhat=f(x)$};
    \draw[-{Stealth}] (D) -- (L);
    \draw[-{Stealth}] (L) -- (f);
    \draw[-{Stealth}] (x) -- (f);
    \draw[-{Stealth}] (f) -- (yh);
  \end{tikzpicture}
  \caption{The supervised learning pipeline. Training happens once, on the left;
  prediction happens many times, on the right. The two are separated in time,
  and confusing them is the source of most leakage bugs (\cref{ch:ch21}).}
  \label{fig:pipeline}
\end{figure}

\section{Supervised Learning}
\label{sec:supervised}

The setting just described --- examples that come with their right answers --- is
\term{supervised learning}. The word \emph{supervised} refers to those answers: a
supervisor has already solved $m$ instances of the problem for us, and the
learner's job is to generalise from them to instances nobody has solved.

That last word carries the weight. A lookup table storing every training pair
reproduces the training answers perfectly and is useless, because it says nothing
about a new $x$. What we want from $f$ is performance on inputs it was not
trained on, and \cref{ch:ch10} is entirely about measuring that honestly.

Supervised problems divide into two kinds, according to what $\Yspace$ looks
like.

\subsection{Regression: predicting a number}

In \term{regression}, the output is a number drawn from a continuous range,
$\Yspace = \R$. \Cref{fig:regression} shows twelve apartments, each plotted by
floor area and sale price, with a straight line fitted through them. To predict
the price of a 1700 square foot apartment, we read the line at $x=17$.

\begin{figure}[htbp]
  \centering
  \begin{tikzpicture}
    \begin{axis}[mlaxis, width=0.78\linewidth, height=5.6cm,
      xlabel={floor area $x$ (hundreds of sq.\ ft)},
      ylabel={price $y$ (\$1000s)},
      xmin=0, xmax=26, ymin=0, ymax=95,
      legend pos=north west]
      \addplot[only marks,mark=*,mark size=1.5pt,MLTeal] coordinates {
        (5,22)(7,26)(9,36)(10,33)(12,44)(13,49)
        (15,52)(16,61)(18,63)(20,72)(22,74)(24,84)};
      \addlegendentry{training examples}
      \addplot[MLRose,very thick,domain=3:25]{3.2984*x+4.3308};
      \addlegendentry{$f(x)=3.2984x+4.3308$}
      \addplot[only marks,mark=diamond*,mark size=3pt,MLAmber,forget plot]
        coordinates {(17,60.404)};
      \draw[dashed,MLAmber,thick] (axis cs:17,0) -- (axis cs:17,60.404)
            -- (axis cs:0,60.404);
    \end{axis}
  \end{tikzpicture}
  \caption{Regression. The fitted line is the least-squares fit to these twelve
  points, computed in \cref{sec:end-to-end}. Reading it at $x=17$ predicts
  $\yhat = 60.4$, or about \$60{,}400.}
  \label{fig:regression}
\end{figure}

The defining feature is that the output space is infinite and ordered. There are
infinitely many prices between \$60{,}000 and \$61{,}000, and a prediction of
\$60{,}400 when the truth is \$61{,}000 is nearly right --- a fact the cost
function will exploit. Other regression problems: forecasting tomorrow's
rainfall, predicting a student's final grade, estimating the remaining life of a
battery, predicting a delivery time.

\subsection{Classification: predicting a category}

In \term{classification}, the output is one of a small, finite set of categories:
$\Yspace = \{0,1\}$ for \term{binary classification}, or $\{1,\dots,K\}$ for
\term{multiclass classification}. The categories need not be numbers at all;
$\{\text{benign}, \text{malignant}\}$ is a perfectly good $\Yspace$, and we
encode it numerically only for the algebra's convenience.

\begin{figure}[htbp]
  \centering
  \begin{tikzpicture}
    \begin{axis}[mlaxis, width=0.78\linewidth, height=5.4cm,
      xlabel={tumour size}, ylabel={patient age},
      xmin=0, xmax=10, ymin=15, ymax=75,
      legend pos=south east]
      \addplot[only marks,mark=o,mark size=2pt,MLTeal,thick] coordinates {
        (1.2,25)(2,32)(1.6,45)(2.8,28)(3.2,38)(2.2,52)(3.8,30)(1.4,58)(4.2,24)};
      \addlegendentry{benign}
      \addplot[only marks,mark=x,mark size=3pt,MLRose,thick] coordinates {
        (6,48)(7.2,55)(5.4,62)(8,45)(6.8,38)(7.8,66)(5.8,52)(8.6,58)};
      \addlegendentry{malignant}
      \addplot[MLBlue,very thick,domain=1.1:7.6,forget plot]{-8*x+76};
    \end{axis}
  \end{tikzpicture}
  \caption{Classification with two features. The learner's output is a
  \term{decision boundary}: points on one side are predicted benign, points on
  the other malignant. \Cref{ch:ch06} derives a principled way to place it.}
  \label{fig:classification}
\end{figure}

\begin{pitfall}{Regression or classification?}
The test is not whether the target is written as a number. It is whether the
\emph{distance} between two outputs is meaningful. Prices of 52 and 53 are nearly
the same, so price is a regression target. If you encode
$\{\text{cat},\text{dog},\text{bird}\}$ as $\{1,2,3\}$, the encoding claims that
a dog is halfway between a cat and a bird, and that a bird is three times a cat.
It is not, and \cref{ch:ch09} shows how to encode categories without smuggling in
an order that does not exist.

The awkward middle case is an ordered category: a star rating from 1 to 5, or a
severity grade. Distances there are partly meaningful, and either treatment can
be defended. Say which you chose and why.
\end{pitfall}

\section{Unsupervised Learning}
\label{sec:unsupervised}

Now remove the labels. We are given only $\{\xii{1},\dots,\xii{m}\}$, with no
$y$ of any kind, and asked to find \emph{structure}. This is
\term{unsupervised learning}, and it is what you do when nobody has solved any
instances for you --- which, outside textbooks, is the common case, because
labels are expensive and data is not.

Three shapes of structure recur, and each gets a chapter later in the book.

\begin{figure}[htbp]
  \centering
  \begin{minipage}[t]{0.32\linewidth}
    \centering
    \begin{tikzpicture}
      \begin{axis}[width=3.9cm,height=3.2cm,axis lines=box,
        xtick=\empty,ytick=\empty,xmin=0,xmax=10,ymin=0,ymax=10,
        title={\scriptsize Clustering}, title style={font=\scriptsize}]
        \addplot[only marks,mark=*,mark size=1.3pt,MLTeal]
          coordinates {(1.5,7.5)(2,8.4)(2.6,7)(1.2,6.6)(2.4,8.9)};
        \addplot[only marks,mark=*,mark size=1.3pt,MLAmber]
          coordinates {(7.6,7.8)(8.4,8.6)(8.8,7.2)(7.2,8.8)};
        \addplot[only marks,mark=*,mark size=1.3pt,MLRose]
          coordinates {(4.6,2.2)(5.4,1.5)(5.8,2.8)(4.2,1.4)(6,1.9)};
      \end{axis}
    \end{tikzpicture}
    {\scriptsize Group similar items.\par}
  \end{minipage}\hfill
  \begin{minipage}[t]{0.32\linewidth}
    \centering
    \begin{tikzpicture}
      \begin{axis}[width=3.9cm,height=3.2cm,axis lines=box,
        xtick=\empty,ytick=\empty,xmin=0,xmax=10,ymin=0,ymax=10,
        title={\scriptsize Anomaly detection}, title style={font=\scriptsize}]
        \addplot[only marks,mark=*,mark size=1.3pt,MLTeal] coordinates {
          (4,4.4)(4.6,5)(5.2,4.2)(3.8,5.4)(5,5.6)(4.4,4)(5.6,4.8)(4.2,5.8)};
        \addplot[only marks,mark=star,mark size=3.4pt,MLRose,thick]
          coordinates {(8.6,1.6)};
        \draw[MLRose,dashed] (axis cs:8.6,1.6) circle[radius=10pt];
      \end{axis}
    \end{tikzpicture}
    {\scriptsize Flag the unusual.\par}
  \end{minipage}\hfill
  \begin{minipage}[t]{0.32\linewidth}
    \centering
    \begin{tikzpicture}
      \begin{axis}[width=3.9cm,height=3.2cm,axis lines=box,
        xtick=\empty,ytick=\empty,xmin=0,xmax=10,ymin=0,ymax=10,
        title={\scriptsize Dimensionality reduction},
        title style={font=\scriptsize}]
        \addplot[only marks,mark=*,mark size=1.3pt,MLTeal] coordinates {
          (1.6,2.4)(2.6,3.2)(3.6,4.6)(4.6,5)(5.6,6.4)(6.6,7)(7.6,8.4)
          (3,3.4)(5,5.8)};
        \addplot[MLBlue,very thick,domain=1:8.4]{1.05*x+0.8};
      \end{axis}
    \end{tikzpicture}
    {\scriptsize Compress many features into few.\par}
  \end{minipage}
  \caption{Three kinds of structure found without labels, treated in
  \cref{ch:ch13,ch:ch14,ch:ch16} respectively.}
  \label{fig:unsupervised}
\end{figure}

\textbf{Clustering} partitions the examples into groups whose members resemble
one another: news articles by topic, customers by buying pattern, genes by
expression profile. Nobody supplies the groups; the algorithm proposes them, and
a human decides afterwards whether they mean anything. \Cref{ch:ch13} develops
$k$-means.

\textbf{Anomaly detection} builds a model of what normal data looks like and
flags examples that do not fit: a fraudulent transaction, an engine about to
fail, a server behaving unlike its peers. \Cref{ch:ch14} does this with a
density estimate.

\textbf{Dimensionality reduction} replaces many features with a few that retain
most of the information, for visualization, for compression, or to make a
downstream model cheaper. \Cref{ch:ch16} derives principal component analysis.

\begin{keyinsight}{The two settings in one line}
Supervised: \emph{here are the right answers; learn to reproduce them on new
inputs.} Unsupervised: \emph{here is the data; tell me something about its
structure.} The first has an objective measure of success on held-out data. The
second usually does not, which is why unsupervised results need more careful
justification, not less.
\end{keyinsight}

\section{Reinforcement Learning}
\label{sec:rl-preview}

A third setting sits alongside these two and is the subject of \cref{ch:ch17}.
In \term{reinforcement learning} there are no labelled examples at all. An
\term{agent} takes \term{actions} in an \term{environment}, the environment
returns a new \term{state} and a scalar \term{reward}, and the agent must learn a
\term{policy} --- a rule for choosing actions --- that maximises reward
accumulated over time.

\begin{figure}[htbp]
  \centering
  \begin{tikzpicture}[every node/.style={font=\scriptsize}]
    \node[draw,rounded corners,fill=MLTeal,text=white,minimum width=2.4cm,
          minimum height=8mm] (ag) {Agent};
    \node[draw,rounded corners,fill=MLBlue,text=white,minimum width=2.4cm,
          minimum height=8mm,below=18mm of ag] (env) {Environment};
    \draw[-{Stealth},thick] (ag.east) to[bend left=55]
      node[midway,right=1pt,align=left] {action $a_t$} (env.east);
    \draw[-{Stealth},thick] (env.west) to[bend left=55]
      node[midway,left=1pt,align=right]
        {state $s_{t+1}$\\ reward $r_{t+1}$} (ag.west);
  \end{tikzpicture}
  \caption{The reinforcement learning loop. No example says which action was
  correct; the only signal is a reward that may arrive many steps later.}
  \label{fig:rl}
\end{figure}

What makes this hard, and separate from supervised learning, is that the reward
is \emph{delayed} and \emph{evaluative}. A chess engine that loses on move 60
receives no instruction as to which of its moves was the mistake, only that the
sequence ended badly. Assigning credit backwards through time is the whole
problem, and \cref{ch:ch17} solves it with the Bellman equation.

Use this rough guide for classifying a problem you meet:

\begin{center}
\small
\renewcommand{\arraystretch}{1.25}
\begin{tabular}{@{}L{4.2cm}L{6.8cm}@{}}
\toprule
\textbf{You have} & \textbf{You are doing}\\
\midrule
inputs with correct outputs & supervised learning
  (regression if $y$ is a number, classification if $y$ is a category)\\
inputs only & unsupervised learning\\
an environment that responds to actions with rewards &
  reinforcement learning\\
\bottomrule
\end{tabular}
\end{center}

\section{A First End-to-End Example}
\label{sec:end-to-end}

Abstractions settle faster when attached to numbers, so here is a complete, tiny
project. Everything in it is developed properly later; the point now is to see
the whole shape at once.

\textbf{The problem.} An estate agency wants to quote a price for an apartment
from its floor area alone. Following \cref{sec:defining}: the task $T$ is
predicting a sale price; the performance measure $P$ is the average squared
error of the prediction; the experience $E$ is twelve recorded sales.

\textbf{The data.} Dataset~D1 of \cref{ch:appD}, plotted in
\cref{fig:regression}. Floor area $x$ is in hundreds of square feet, price $y$
in thousands of dollars, and $m=12$.

\textbf{The hypothesis space.} Straight lines, $f(x)=wx+b$. This is a strong
assumption --- prices are not exactly linear in area --- and \cref{ch:ch05}
relaxes it.

\textbf{The cost.} Squared error, $J(w,b)=\frac{1}{2m}\sum_i (f(\xii{i})-\yii{i})^2$.

\textbf{The optimizer.} Gradient descent, which \cref{ch:ch03} derives. Started
at $w=b=0$ with a step size of $0.005$, it produces the trajectory in
\cref{fig:gd-trace}: the cost falls from $1499.7$ to $3.7$ within ten
iterations, then spends thousands more inching to its final value.

\begin{workedexample}{Fitting, scoring, and testing the apartment model}
\textbf{1. The fitted model.} Minimising $J$ over the twelve points gives
\[
  w = 3.2984, \qquad b = 4.3308, \qquad J(w,b) = 2.4845 .
\]
Interpretation: each additional hundred square feet is worth about \$3{,}300,
and the line meets the axis at \$4{,}330 --- a number with no physical meaning,
since an apartment of zero area does not sell for \$4{,}330. Intercepts often
lack an interpretation; that is not an error.

\textbf{2. A prediction.} For $x = 17$,
$\yhat = 3.2984(17) + 4.3308 = 60.4040$, so about \$60{,}400.

\textbf{3. Is that any good?} The number $J=2.4845$ means nothing alone. Compare
it with the simplest model that ignores the input entirely: always predict the
mean price, $\bar y = 51.3333$. That constant model has $J = 182.1111$. The
fitted line reduces the cost by a factor of $73$, and its root mean squared
error is $2.23$ against the baseline's $19.08$. \emph{Now} we know something.

\textbf{4. Do other lines do better?} No, and the cost says so:
\begin{center}
\small
\begin{tabular}{@{}lrr@{}}
\toprule
model & $J(w,b)$ & RMSE\\
\midrule
$f(x)=3.2984x+4.3308$ \quad(the fit) & $2.4845$  & $2.2291$\\
$f(x)=3.0x+10.0$                     & $4.9583$  & $3.1491$\\
$f(x)=3.6x+3.0$                      & $8.3867$  & $4.0955$\\
$f(x)=4.0x$                          & $26.6667$ & $7.3030$\\
$f(x)=0$                             & $1499.6667$ & $54.7662$\\
\bottomrule
\end{tabular}
\end{center}
Note the third line: a plausible-looking fit, eyeballed onto the same scatter,
costs more than three times the optimum. Eyeballing is not fitting.

\textbf{5. The honest test.} All of the above is measured on the same twelve
points the model was fitted to, which flatters it. Refit using only the nine
smallest apartments and predict the three largest, which the model has never
seen:
\[
  w = 3.3626,\qquad b = 3.6584,
\]
\begin{center}
\small
\begin{tabular}{@{}rrrr@{}}
\toprule
$x$ & actual $y$ & predicted $\yhat$ & error\\
\midrule
$20$ & $72$ & $70.91$ & $-1.09$\\
$22$ & $74$ & $77.64$ & $+3.64$\\
$24$ & $84$ & $84.36$ & $+0.36$\\
\bottomrule
\end{tabular}
\end{center}
Training RMSE $2.30$, test RMSE $2.20$. The model does about as well on
apartments it has never seen as on the ones it was fitted to, which is the
outcome we want and, as \cref{ch:ch07} shows, not the outcome we always get.
\end{workedexample}

\begin{figure}[htbp]
  \centering
  \begin{tikzpicture}
    \begin{axis}[mlaxis, width=0.78\linewidth, height=5.2cm,
      xmode=log, ymode=log,
      xlabel={iteration}, ylabel={cost $J(w,b)$},
      xmin=1, xmax=20000, ymin=2, ymax=100]
      \addplot[MLBlue,thick,mark=*,mark size=0.9pt] coordinates {
        (1,54.7135)(2,5.4081)(3,3.7241)(4,3.6651)(6,3.6598)(7,3.6581)
        (10,3.6532)(13,3.6484)(18,3.6403)(24,3.6306)(32,3.6179)(42,3.6022)
        (56,3.5806)(75,3.5520)(100,3.5154)(133,3.4691)(178,3.4092)
        (237,3.3362)(316,3.2474)(422,3.1426)(562,3.0259)(750,2.9011)
        (1000,2.7785)(1334,2.6691)(1778,2.5839)(2371,2.5280)(3162,2.4989)
        (4217,2.4878)(5623,2.4850)(7499,2.4845)(10000,2.4845)(13335,2.4845)
        (17783,2.4845)};
      \addlegendentry{gradient descent, $\alpha=0.005$}
      \addplot[MLRose,dashed,thick,domain=1:20000,forget plot]{2.4845};
    \end{axis}
  \end{tikzpicture}
  \caption{Cost against iteration, on log axes. The dashed line is the minimum,
  $J=2.4845$. Most of the improvement happens in the first ten steps; the long
  tail is the slow drift of the intercept, and \cref{ch:ch05} explains why it is
  slow and how feature scaling fixes it.}
  \label{fig:gd-trace}
\end{figure}

\begin{keyinsight}{Two numbers, never one}
A model's error is meaningless in isolation. Always report it against a baseline
--- the constant prediction, the existing rule, the human --- and always measure
it on data the model has not seen. An error rate with neither comparison nor a
held-out set is a claim about memory, not about learning.
\end{keyinsight}

\begin{labbox}{Lab 1 --- Environment, data, and a first model}
Set up Python, NumPy and Jupyter; load dataset D1; reproduce
\cref{fig:regression}; and write a function \code{cost(x, y, w, b)} that returns
$J$. Check it against the table in the worked example: your function must return
$2.4845$ at $(3.2984, 4.3308)$ and $1499.6667$ at $(0,0)$. Two near misses tell
you exactly what went wrong: $4.9690$ means you omitted the $\tfrac12$, and
$29.8141$ means you omitted the $\tfrac{1}{m}$.
\end{labbox}

\section{What Machine Learning Is Not}
\label{sec:what-its-not}

The failures worth learning about early are not algorithmic. They are failures of
judgement about whether learning applies at all.

\subsection{It is not a substitute for a known rule}

If the rule can be written, write it. A learned model that computes a payroll
deduction will be right \emph{almost} always, will fail unpredictably on rare
inputs, cannot be audited line by line, and requires a data pipeline to keep
working. A dozen lines of arithmetic have none of those properties. Use learning
where specification by example is the only specification available.

\subsection{It does not discover causes}

A model finds patterns that predict, and prediction is not causation. A model
trained on hospital records may learn that patients who receive a particular
drug die more often, because the drug is given to the sickest patients. The
correlation is real and useful for prediction --- and acting on it as though it
were a cause would kill people.

\begin{pitfall}{The direction of the arrow}
Ask of any model you are tempted to act on: am I \emph{predicting} an outcome, or
am I \emph{intervening} to change it? Predictive accuracy licenses the first and
says nothing about the second. Every feature that works because it is a
\emph{consequence} of the target rather than a cause of it --- a treatment code,
a follow-up appointment, a flag written after the fact --- will predict
beautifully in training and fail the moment the model is used to make decisions.
\Cref{ch:ch21} treats this as a data-leakage problem, which is what it is.
\end{pitfall}

\subsection{It cannot outrun its data}

The model inherits every property of the training set, including the ones nobody
intended. If the recorded past decisions were made by people, the model learns
their judgement, including its errors and its prejudices. If some group of inputs
is rare in the data, the model will be worse for that group, and the overall
accuracy figure will hide it, because that group is rare in the test set too.

\begin{ethicsnote}{Whose examples are they?}
When \textcite{buolamwini2018} evaluated three commercial gender-classification
systems, every one of them reported strong overall accuracy. Disaggregated, the
error rate for lighter-skinned men was under $1\%$, while for darker-skinned
women it exceeded $30\%$ in the worst system. Nothing was wrong with the
algorithms in any sense a course on optimization would recognise. The training
sets simply contained far more of some faces than others, and a single averaged
accuracy figure is arithmetically incapable of revealing that.

Two habits follow, and they cost almost nothing to adopt now: know the
composition of your training data, and report performance for each group your
system will be applied to, not only the average. \Cref{ch:ch11} makes both
routine.
\end{ethicsnote}

\subsection{It is not finished when it is accurate}

A model that scores well in a notebook is roughly a third of a working system.
The rest --- serving predictions within a latency budget, noticing when incoming
data stops resembling training data, retraining, and rolling back a bad release
--- is the subject of \cref{ch:ch20,ch:ch21,ch:ch22}, and is where most
industrial machine learning projects actually fail.

\section{A Map of This Book}
\label{sec:map}

\Cref{tab:map} shows the route. Parts~I and~II are the load-bearing ones: they
develop cost functions, gradient descent, overfitting, regularization and
evaluation on models simple enough to draw on a page and check by hand. Almost
everything later is a variation on machinery introduced there.

\begin{table}[htbp]
  \centering
  \small
  \renewcommand{\arraystretch}{1.25}
  \begin{tabular}{@{}L{1.1cm}L{4.4cm}L{5.6cm}@{}}
    \toprule
    \textbf{Part} & \textbf{Subject} & \textbf{The question it answers}\\
    \midrule
    I    & Linear regression &
      How does fitting work, mechanically?\\
    II   & Classification and regularization &
      How do we predict categories, and why does a model that fits perfectly
      often predict badly?\\
    III  & Neural networks &
      What if the right features are themselves unknown?\\
    IV   & Diagnosis and improvement &
      The model is not good enough. What should I do next?\\
    V    & Unsupervised learning &
      What can be learned with no labels at all?\\
    VI   & Recommenders and PCA &
      How do we learn about users and items together, and how do we compress
      features?\\
    VII  & Reinforcement learning &
      How does an agent learn from delayed reward?\\
    VIII & Transformers and language models &
      Where does the current generation of language models come from?\\
    IX   & Machine learning systems &
      It works in a notebook. Now what?\\
    \bottomrule
  \end{tabular}
  \caption{The nine parts and what each is for.}
  \label{tab:map}
\end{table}

\section*{Chapter Summary}
\addcontentsline{toc}{section}{Chapter Summary}

\begin{itemize}
  \item Machine learning inverts classical programming: examples of correct
        answers go in, and rules come out. It is the right tool when a rule
        exists, cannot be stated, and examples of its output are plentiful.
  \item Mitchell's definition makes a project concrete by forcing you to name
        the task $T$, the performance measure $P$, and the experience $E$.
        Accuracy is frequently the wrong $P$.
  \item Every supervised learner has four parts: a training set, a hypothesis
        space $\Hspace$ chosen by the designer, a cost function $J$ measuring
        fit, and an optimizer that searches $\Hspace$ for the cheapest
        hypothesis.
  \item Supervised learning uses labelled examples, and splits into regression
        (continuous output) and classification (categorical output). The test
        is whether distance between two outputs is meaningful.
  \item Unsupervised learning finds structure in data that carries no labels
        at all: clustering, anomaly detection, and dimensionality reduction.
  \item Reinforcement learning learns a policy from delayed, evaluative reward.
  \item A cost is meaningless in isolation. Compare it against a baseline, and
        measure it on data the model has not seen.
  \item Machine learning does not discover causes, cannot outrun its training
        data, and is not finished when it is accurate.
\end{itemize}

\section*{Exercises}
\addcontentsline{toc}{section}{Exercises}

\exercisesection{Short Questions}

\sq Give one task for which machine learning is clearly the wrong tool, and one
for which it is clearly the right one. Justify each in a single sentence using
the criteria of \cref{sec:applications}.

\sq State Mitchell's definition of learning, and identify $T$, $P$ and $E$ for a
system that learns to route a delivery van.

\sq Define the hypothesis space. Why does the learning algorithm not simply
search over all functions from $\Xspace$ to $\Yspace$?

\sq What does the superscript in $\xii{3}$ denote? Write an expression in which
both $\xii{3}$ and $x^3$ appear, and say what each means.

\sq Explain, in two sentences, the difference between regression and
classification, without using the words \emph{continuous} or \emph{discrete}.

\sq A dataset has $m = 400$ examples and each example has $12$ measurements.
What are $m$ and $n$? What is the shape of the design matrix $\Xmat$?

\sq Name the three shapes of structure that unsupervised learning looks for, and
give an application of each not mentioned in this chapter.

\sq Why is a lookup table that stores every training pair a poor model, despite
achieving zero error on the training set?

\sq What is a baseline, and why is a reported error rate without one close to
useless?

\sq In one sentence each, say why machine learning cannot establish causation
and why it cannot exceed the quality of its training data.

\exercisesection{Analytical Problems}

\ap For each of the following, state whether it is supervised regression,
supervised classification, unsupervised, or reinforcement learning, and say what
$x$ and $y$ (if any) are.
\begin{enumerate}[label=(\alph*)]
  \item From $10{,}000$ patient records carrying a diabetic / non-diabetic flag,
        predict whether a new patient is diabetic.
  \item Given browsing logs for $50{,}000$ users with no annotation, split them
        into groups with similar interests.
  \item Given yesterday's electricity load curve, predict today's peak demand in
        megawatts.
  \item Given millions of card transactions, none marked fraudulent, find the
        ones worth investigating.
  \item Train a robot arm to place a component on a circuit board, where the only
        feedback is whether the finished board passes testing.
  \item Given $60{,}000$ images of handwritten digits, each labelled $0$--$9$,
        read the digit in a new image.
\end{enumerate}

\ap A screening test is applied to a population in which $2\%$ of people have the
condition. A program that always answers ``healthy'' is proposed.
\begin{enumerate}[label=(\alph*)]
  \item What accuracy does it achieve?
  \item Describe, in words, a performance measure $P$ under which this program
        scores badly.
  \item Which of the two errors --- missing a sick patient, or alarming a healthy
        one --- would you weight more heavily, and what does your answer imply
        about $P$?
\end{enumerate}

\ap Using the twelve apartments of \cref{fig:regression}, compute $J(w,b)$ by
hand for $w = 3$, $b = 5$. Show the residual for each of the first four
examples, then the total. Is your answer larger or smaller than the $2.4845$
achieved by the fitted line? Explain why it must be.

\ap The constant model $f(x)=c$ has one parameter. Show that the value of $c$
minimising $\frac{1}{2m}\sum_i (c - \yii{i})^2$ is $c = \bar y$, the mean of the
targets. (Differentiate with respect to $c$ and set the derivative to zero.)
Verify against the value $\bar y = 51.3333$ used in \cref{sec:end-to-end}.

\ap Suppose you encode $\{\text{low}, \text{medium}, \text{high}\}$ as
$\{1,2,3\}$ and fit a regression model. Give a concrete prediction the model
might make that reveals the encoding to be inappropriate, and describe an input
for which the model's output would be uninterpretable.

\ap A colleague reports that their model achieves $94\%$ accuracy and proposes
deployment. List four questions you would ask before agreeing, and say what a
bad answer to each would look like.

\exercisesection{Design Problems}

\dpr A university wants to predict which first-year students are at risk of
failing, so that tutors can offer support early. Specify $T$, $P$ and $E$; name
the features you would use and at least one you would deliberately refuse to use,
with your reason; state whether this is regression or classification; and
describe how you would measure success six months after deployment. Then argue
the other side: give the strongest case that this system should not be built.

\dpr You are asked to build a model that predicts crop yield from weather and
soil measurements, for smallholder farmers who will decide what to plant.
Labelled data exists for three of the twelve districts in the region. Describe
how you would proceed, what you would measure, and what you would say to a farmer
in one of the other nine districts about how much to trust the prediction.

\dpr Choose an application from \cref{tab:applications} and write its
$\langle T, P, E\rangle$, then identify the baseline that any learned model must
beat. Estimate --- with reasoning, not a guess --- how much better than that
baseline the model must be for the system to be worth building.

\dpr A team proposes replacing a rule-based fraud check with a learned model.
The existing rules are auditable and are cited in regulatory filings. Write half
a page for the decision-maker setting out what is gained, what is lost, and what
you would need to see before recommending the change.

\exercisesection{Programming Exercises}

\pe Enter dataset D1 as two NumPy arrays and reproduce \cref{fig:regression},
including the dashed prediction lines at $x = 17$.

\pe Write \code{cost(x, y, w, b)} returning $J(w,b)$, with no Python loop over
the examples. Reproduce every row of the cost table in the worked example. Then
write a second version using an explicit loop and assert that the two agree to
within $10^{-12}$.

\pe Scan $w$ over the range $[2, 5]$ in steps of $0.01$ with $b$ fixed at
$4.3308$, plot $J$ against $w$, and read off the minimising $w$ from the plot.
How close is it to $3.2984$? What does the shape of the curve tell you about how
hard this optimization is?

\pe Split D1 into the first nine examples and the last three. Fit on the nine
using the closed-form solution
$w = \sum_i (\xii{i}-\bar x)(\yii{i}-\bar y) \big/ \sum_i (\xii{i}-\bar x)^2$
and $b = \bar y - w\bar x$, and reproduce the held-out table in the worked
example. Then split the other way --- train on the three largest, test on the
nine smallest --- and comment on the difference.

\exercisesection{Self-Test}

\qz A program's performance on a task improves as it processes more data. Under
Mitchell's definition, the data is: (a)~$T$; (b)~$P$; (c)~$E$; (d)~$\Hspace$.

\qz Predicting the number of days until a machine fails is: (a)~classification;
(b)~regression; (c)~clustering; (d)~reinforcement learning.

\qz Grouping unlabelled documents by topic is: (a)~supervised;
(b)~unsupervised; (c)~reinforcement; (d)~none of these.

\qz In the notation of this book, $\xjj{4}{2}$ means: (a)~$x$ raised to a power;
(b)~feature 2 of example 4; (c)~example 2 of feature 4; (d)~the fourth
derivative.

\qz A model has zero error on its training set and high error on new data. The
most likely explanation is: (a)~the cost function is wrong; (b)~the hypothesis
space is too small; (c)~the model has memorised the training set; (d)~the
learning rate was too large.

\qz A classifier for a condition affecting $1$ in $1000$ people reports $99.9\%$
accuracy. The most useful next question is: (a)~what was the learning rate;
(b)~how does it compare with always predicting ``healthy''; (c)~how many
features were used; (d)~how long did training take.

\qz Which is \emph{not} one of the four components of a supervised learning
algorithm? (a)~a hypothesis space; (b)~a cost function; (c)~a test set;
(d)~an optimizer.

\qz A model trained on past hiring decisions reproduces a historical bias against
a group. The best description is: (a)~a bug in the optimizer; (b)~overfitting;
(c)~the model faithfully learning a property of its training data;
(d)~an error in the cost function.

\furtherreading

\textcite{james2021} open with the same distinction between supervised and
unsupervised learning; read §1 for context and §2.1--2.2 for the statistical
framing of prediction, including the difference between prediction and
inference, which is the formal version of \cref{sec:what-its-not}'s warning
about causation. \textcite{bishop2006} §1.1 covers the same ground through the
polynomial curve-fitting example, and is worth reading now and again after
\cref{ch:ch07}, because it introduces overfitting in the first ten pages.

\bigskip
\noindent\emph{Next:} \Cref{ch:ch02} takes the fourth component apart. We fix the
hypothesis space to straight lines and ask, precisely, what makes one line better
than another --- which turns out to be the entire content of the cost function.
